% 初出一覧（原书 p.166）——Lead
\chapter*{初出一览}
\addcontentsline{toc}{chapter}{初出一览（原书 p.166）}
\markboth{初出一览}{初出一览}

本书各章取自『数学セミナー』2006 年 10 月号至 2007 年 8 月号的连载（第 3 章除外），
成书时作了若干增补与修改。各章出处与执笔者如下（括号内为执笔当时的所属）。

\begin{description}[leftmargin=1.5em,itemsep=6pt,style=nextline]
  \item[第 1 章\quad 什么是镜像对称] 『数学セミナー』2006 年 10 月号\\
    江口徹（京都大学基础物理学研究所）
  \item[第 2 章\quad Gromov--Witten 不变量入门] 『数学セミナー』2006 年 11 月号\\
    深谷賢治（京都大学研究生院理学研究科）
  \item[第 3 章\quad 正多面体与镜像对称] 新写\\
    高橋篤史（大阪大学研究生院理学研究科）
  \item[第 4 章\quad Calabi--Yau 流形与代数几何] 『数学セミナー』2006 年 12 月号\\
    並河良典（京都大学研究生院理学研究科）
  \item[第 5 章\quad 奇点与拓扑场论] 『数学セミナー』2007 年 1 月号（原题「ミラー対称性と特異点 I」）\\
    高橋篤史（大阪大学研究生院理学研究科）
  \item[第 6 章\quad 镜像对称与共形场论] 『数学セミナー』2007 年 2 月号\\
    細野忍（东京大学研究生院数理科学研究科）
  \item[第 7 章\quad 环簇几何与 Calabi--Yau 流形] 『数学セミナー』2007 年 3 月号（原题「ミラー対称性とトーリック幾何学」）\\
    今野宏（东京大学研究生院数理科学研究科）
  \item[第 8 章\quad Floer 同调与镜像对称] 『数学セミナー』2007 年 4 月号\\
    深谷賢治（京都大学研究生院理学研究科）
  \item[第 9 章\quad 镜像对称与 D-模] 『数学セミナー』2007 年 5 月号\\
    Martin Guest（首都大学东京研究生院理工学研究科）、乙藤隆史（日本大学工学部）
  \item[第 10 章\quad 镜像对称与同调代数] 『数学セミナー』2007 年 6 月号\\
    深谷賢治（京都大学研究生院理学研究科）
  \item[第 11 章\quad 奇点与镜像对称] 『数学セミナー』2007 年 7 月号（原题「ミラー対称性と特異点 II」）\\
    高橋篤史（大阪大学研究生院理学研究科）
  \item[第 12 章\quad 超越镜像对称] 『数学セミナー』2007 年 8 月号\\
    中島啓（京都大学数理解析研究所）
\end{description}
